% !TeX program = xelatex
% ============================================================
% ro/ro.tex — Руководство оператора
% ГОСТ 19.505-79 «Руководство оператора. Требования к содержанию
% и оформлению» (ЕСПД).
% ============================================================
\documentclass[12pt,a4paper]{extarticle}

% 1. Подключение общей преамбулы (пакеты)
\input{../common/preamble}

% 2. Определение переменных для конкретного документа.
%    Базовый код RU.… приходит из \hseDocCodeBase (common/meta.tex);
%    здесь задаётся только суффикс типа документа по ГОСТ 19.103-77.
\newcommand{\docname}{Руководство оператора}
\newcommand{\doccode}{\hseDocCodeBase\ 34 01-1}
\newcommand{\doccodelu}{\hseDocCodeBase\ 34 01-1-ЛУ}

% 3. Подключение общих команд (проект, студент, руководитель)
\input{../common/commands}

% 4. Подключение стилей (колонтитулы, заголовки)
\input{../common/styles}

\begin{document}
\sloppy

% 5. Генерация титульных листов по шаблону
\input{../common/title_template}


% ============================================================
% АННОТАЦИЯ
% ============================================================
\section*{АННОТАЦИЯ}
\addcontentsline{toc}{section}{АННОТАЦИЯ}

% Подсказка: уточните формулировки под вашу систему (что именно запускает
% и контролирует оператор).
Настоящий документ «\projectname. Руководство оператора» содержит сведения, необходимые оператору для запуска программы, контроля состояния системы и выполнения типовых эксплуатационных проверок в программном комплексе.

В документе описаны назначение программы, условия ее выполнения, последовательность действий оператора при работе с системой, а также сообщения, выдаваемые программой в ходе выполнения пользовательских и эксплуатационных сценариев.

Руководство предназначено для операторов, администраторов и инженеров эксплуатации, выполняющих первичный контроль работоспособности, проверку объектов системы, просмотр журналов, метрик и результатов выполнения операций.

Документ разработан в соответствии с требованиями ГОСТ 19.505-79\cite{gost19505}.

\clearpage

% ============================================================
% СПИСОК ИСПОЛЬЗУЕМЫХ ГОСТ
% ============================================================
\noindent
Настоящий документ разработан в соответствии с требованиями:

\begin{enumerate}[label=\arabic*), leftmargin=1.25cm]
\item ГОСТ 19.101-77 Виды программ и~программных документов\cite{gost19101}.
\item ГОСТ 19.103-77 Обозначения программ и~программных документов\cite{gost19103}.
\item ГОСТ 19.104-78 Основные надписи\cite{gost19104}.
\item ГОСТ 19.105-78 Общие требования к~программным документам\cite{gost19105}.
\item ГОСТ 19.106-78 Требования к~программным документам, выполненным печатным способом\cite{gost19106}.
\item ГОСТ 19.505-79 Руководство оператора. Требования к~содержанию и~оформлению\cite{gost19505}.
\end{enumerate}

Изменения к~данному документу оформляются согласно ГОСТ\,19.604-78\cite{gost19604}.

\clearpage

% ============================================================
% СОДЕРЖАНИЕ
% ============================================================
\hypersetup{linkcolor=black}
\tableofcontents
\hypersetup{linkcolor=blue}
\thispagestyle{fancy}
\clearpage

% ============================================================
% 1. НАЗНАЧЕНИЕ ПРОГРАММЫ
% ============================================================
\section{НАЗНАЧЕНИЕ ПРОГРАММЫ}

\subsection{Назначение программы}

% Подсказка: опишите, для чего предназначена программа, из каких основных
% частей (сервисов, модулей) она состоит и через какой интерфейс
% оператор с ней взаимодействует.
\begin{hseExample}
Программа «\projectname» предназначена для \hseFill{кратко назначение системы}. Реализованная часть системы включает \hseFill{перечислите основные сервисы или модули}. Через \hseFill{основной модуль или подсистему} выполняется \hseFill{перечислите ключевые операции системы}.

Оператор использует программу для контроля состояния системы, выполнения типовых проверочных действий, проверки результата операций и первичного анализа отказов. Основное взаимодействие выполняется через \hseFill{интерфейс взаимодействия, например REST API}. В локальном контуре отдельные сервисы могут проверяться напрямую согласно инструкциям репозитория проекта.
\end{hseExample}

\subsection{Основные эксплуатационные функции}

\begin{hseExample}
С точки зрения оператора программа выполняет следующие функции:

% Подсказка: первые пункты замените на операции над объектами и сценарии
% вашей системы; последние два пункта — универсальные, их можно оставить.
\begin{itemize}[leftmargin=1.25cm]
    \item создание, просмотр, изменение и удаление \hseFill{основные объекты системы} через доступные API;
    \item выполнение пользовательских сценариев \hseFill{перечислите ключевые сценарии};
    \item контроль состояния сервисов по журналам, метрикам, трассировкам и сообщениям об ошибках;
    \item проверка готовности сервисов и инфраструктурных зависимостей.
\end{itemize}

Настоящее руководство описывает только функции, подтвержденные реализацией и развертываемыми конфигурациями.
\end{hseExample}

\subsection{Режимы работы}

\begin{hseExample}
Программа может использоваться в следующих режимах:

% Подсказка: скорректируйте перечень режимов под ваш проект.
\begin{enumerate}[label=\arabic*), leftmargin=1.25cm]
    \item \textbf{Локальный режим}: сервисы и инфраструктурные зависимости запускаются в контейнерном окружении разработки. Режим применяется для первичной проверки и демонстрации сценариев.
    \item \textbf{Стендовый режим}: сервисы развернуты в \hseFill{среда развертывания, например Kubernetes}, доступ к API и средствам наблюдаемости предоставляется через выделенные адреса стенда.
    \item \textbf{Приемочный режим}: оператор выполняет контрольные действия по программе и методике испытаний, фиксирует состояние сервисов и результаты проверок.
\end{enumerate}
\end{hseExample}

\clearpage

% ============================================================
% 2. УСЛОВИЯ ВЫПОЛНЕНИЯ ПРОГРАММЫ
% ============================================================
\section{УСЛОВИЯ ВЫПОЛНЕНИЯ ПРОГРАММЫ}

\subsection{Требования к техническим средствам}

\begin{hseExample}
Для работы оператора с программой требуется рабочая станция со следующими характеристиками:

% Подсказка: укажите минимальные требования к рабочей станции оператора.
\begin{itemize}[leftmargin=1.25cm]
    \item процессор: не менее \hseFill{число ядер} ядер x86\_64;
    \item оперативная память: не менее \hseFill{объем памяти} ГБ;
    \item дисковое пространство: не менее \hseFill{объем диска} ГБ свободного пространства для журналов, временных файлов и экспортируемых отчетов;
    \item сетевой доступ к \hseFill{стенд, репозиторий, средства наблюдаемости}.
\end{itemize}

При выполнении локального запуска или диагностики контейнеров рекомендуется рабочая станция с \hseFill{рекомендуемые характеристики станции} и установленной контейнерной средой выполнения.
\end{hseExample}

\subsection{Требования к программным средствам}

\begin{hseExample}
Для выполнения операторских действий требуются следующие программные средства:

% Подсказка: перечислите ПО, необходимое оператору (браузер, клиенты
% запросов, контейнерная среда, утилиты диагностики).
\begin{itemize}[leftmargin=1.25cm]
    \item современный браузер для доступа к \hseFill{веб-интерфейсы и документация стенда};
    \item клиент HTTP-запросов: \hseFill{например curl, Postman} или встроенный интерфейс OpenAPI;
    \item доступ к репозиторию \hseFill{укажите git-хостинг} и артефактам проекта, если оператору поручена проверка CI/CD;
    \item \hseFill{контейнерная среда, например Docker} при локальной проверке сервисов;
    \item средства просмотра журналов, метрик и трассировок: \hseFill{укажите инструменты наблюдаемости}.
\end{itemize}

Доступ к исходному коду и эксплуатационным конфигурациям предоставляется через закрытый репозиторий \hseFill{укажите git-хостинг проекта}. Репозиторий закрыт, доступ выдается по запросу ответственным лицам проекта.
\end{hseExample}

\subsection{Требования к доступу}

\begin{hseExample}
Перед началом работы оператор должен получить:

\begin{itemize}[leftmargin=1.25cm]
    \item учетную запись пользователя программы «\projectname» или тестовые учетные данные стенда;
    \item набор прав, достаточный для проверяемых действий: \hseFill{перечислите роли и уровни доступа};
    \item адрес \hseFill{точка входа, например API-шлюз} и, при необходимости, адреса сервисов наблюдаемости;
    \item сведения о тестовом контуре: окружение, версия развертывания, набор доступных сервисов и ограничения доступа;
    \item секреты, токены или ключи доступа только в том случае, если они необходимы для порученных эксплуатационных действий.
\end{itemize}

Оператор не должен хранить пароли, токены и ключи в открытом виде в репозитории, документах, сообщениях или журналах. При передаче диагностических материалов чувствительные данные должны быть удалены или замаскированы.
\end{hseExample}

\subsection{Проверка готовности перед началом работы}

\begin{hseExample}
До выполнения основных действий оператор проверяет готовность программы:

\begin{enumerate}[label=\arabic*), leftmargin=1.25cm]
    \item убедиться, что выбран правильный контур: локальный, стендовый или приемочный;
    \item проверить доступность \hseFill{точка входа, например API-шлюз} по адресу стенда или прямой адрес проверяемого сервиса в локальном контуре;
    \item проверить доступность страницы OpenAPI, если она включена для сервиса, или контрольной точки проверки готовности;
    \item убедиться, что основные сервисы имеют состояние готовности;
    \item открыть панель \hseFill{средство мониторинга} и проверить поступление метрик;
    \item при необходимости проверить журналы сервисов и ошибки приложения в \hseFill{средства сбора журналов и ошибок}.
\end{enumerate}

% Подсказка: замените образцы строк таблицы объектами вашей системы
% (точка входа, сервисы, зависимости, средства наблюдаемости).
% Таблица набрана через \captionof (не \begin{table}[H]): плавающие
% окружения внутри жёлтого блока-образца недопустимы.
\begin{center}
\small
\captionof{table}{Контрольные признаки готовности программы}
\label{tab:ro_readiness}
\begin{tabularx}{\textwidth}{|>{\raggedright\arraybackslash}p{0.28\textwidth}|X|X|}
\hline
\textbf{Проверяемый объект} & \textbf{Ожидаемый результат} & \textbf{Действие при отклонении} \\
\hline
\hseFill{точка входа, например API-шлюз} & Доступен по адресу контура, возвращает ответ на запрос проверки готовности & Проверить маршрут, DNS, TLS-сертификат и состояние компонента \\
\hline
Сервисы \hseFill{перечень сервисов} & Сервисы имеют состояние готовности, нет циклических перезапусков & Проверить логи, переменные окружения и доступность зависимостей \\
\hline
\hseFill{зависимости, например СУБД и брокер} & Зависимости доступны для сервисов, ошибки подключения отсутствуют & Проверить состояние контейнеров или рабочих нагрузок \\
\hline
\hseFill{средства наблюдаемости} & Метрики, журналы и трассировки поступают в контур наблюдаемости & Проверить агентов сбора, источники данных и сетевые политики \\
\hline
\end{tabularx}
\end{center}
\end{hseExample}

\clearpage

% ============================================================
% 3. ВЫПОЛНЕНИЕ ПРОГРАММЫ
% ============================================================
\section{ВЫПОЛНЕНИЕ ПРОГРАММЫ}

\subsection{Общий порядок работы оператора}

\begin{hseExample}
Работа оператора выполняется в следующей последовательности:

\begin{enumerate}[label=\arabic*), leftmargin=1.25cm]
    \item подготовить рабочее место и проверить наличие доступа к стенду;
    \item проверить готовность точки входа в систему, а также состояние сервисов и инфраструктурных зависимостей;
    \item выполнить вход в систему или получить токен доступа через согласованный сценарий аутентификации;
    \item выполнить требуемое операторское действие: просмотр, создание, изменение или проверку объекта;
    \item проверить результат операции по ответу API, журналам, событиям и метрикам;
    \item зафиксировать результат проверки или передать инцидент ответственному инженеру при обнаружении отказа.
\end{enumerate}
\end{hseExample}

\subsection{Запуск программы в локальном контуре}

\begin{hseExample}
Локальный запуск применяется для демонстрационных и проверочных действий. Перед запуском оператор должен убедиться, что установлена \hseFill{контейнерная среда, например Docker} и получен доступ к репозиторию проекта.

% Подсказка: замените плейсхолдеры на реальные команды вашего проекта,
% оформляя их через \texttt{...}.
\begin{enumerate}[label=\arabic*), leftmargin=1.25cm]
    \item получить актуальную версию репозитория;
    \item подготовить файл переменных окружения согласно инструкции проекта;
    \item запустить инфраструктурные зависимости командой \hseFill{команда запуска инфраструктуры};
    \item запустить сервисы командой \hseFill{команда запуска сервисов};
    \item проверить запуск сервисов командой \hseFill{команда проверки};
    \item при необходимости остановить локальное окружение командой \hseFill{команда остановки}.
\end{enumerate}
\end{hseExample}

\subsection{Запуск программы в стендовом контуре}

\begin{hseExample}
В стендовом контуре программа развернута в \hseFill{среда развертывания, например Kubernetes}. Оператор не выполняет разработку исходного кода, а контролирует состояние уже развернутых компонентов:

\begin{enumerate}[label=\arabic*), leftmargin=1.25cm]
    \item проверить состояние \hseFill{пространство имен или группа ресурсов};
    \item убедиться, что все требуемые развертывания готовы;
    \item проверить сервисы и маршруты \hseFill{точка входа и ключевые маршруты};
    \item открыть \hseFill{средство мониторинга} и убедиться, что метрики поступают;
    \item выполнить контрольный пользовательский или административный сценарий.
\end{enumerate}

% Подсказка: при наличии прав эксплуатации приведите команды диагностики
% (например, kubectl get pods, kubectl logs) в окружении verbatim.
При наличии прав эксплуатации допускается использовать команды диагностики \hseFill{перечислите команды диагностики}.
\end{hseExample}

\subsection{Вход оператора в систему}

\begin{hseExample}
Для выполнения действий через API оператор должен получить токен доступа. Последовательность входа зависит от включенного сценария аутентификации:

\begin{enumerate}[label=\arabic*), leftmargin=1.25cm]
    \item открыть страницу авторизации или выполнить HTTP-запрос к сервису аутентификации;
    \item указать идентификатор пользователя и пароль;
    \item при необходимости подтвердить вход одноразовым кодом или дополнительным фактором;
    \item получить токен доступа и использовать его в заголовке \texttt{Authorization};
    \item убедиться, что запрос к защищенному API возвращает успешный ответ.
\end{enumerate}

Токен доступа используется только в пределах согласованного тестового или стендового контура. Перед публикацией диагностических материалов токены должны быть удалены.
\end{hseExample}

\subsection{Управление организациями и участниками}

% Подсказка: если в вашей системе нет организаций и участников —
% замените подраздел на управление основными объектами вашей системы.
\begin{hseExample}
Для работы с организационной структурой оператор выполняет следующие действия:

\begin{enumerate}[label=\arabic*), leftmargin=1.25cm]
    \item создать или выбрать организацию;
    \item проверить доступность организации по списку организаций текущего пользователя;
    \item сформировать приглашение для нового участника;
    \item проверить статус приглашения по ответу API;
    \item при необходимости принять приглашение и удалить участника через предусмотренные конечные точки API.
\end{enumerate}

Ожидаемый результат: созданные организации отображаются через API, принятие приглашения возвращает запись участника, а удаление участника завершается успешным HTTP-ответом.
\end{hseExample}

\subsection{Управление топологией инфраструктуры}

% Подсказка: перечислите сущности вашей предметной области и порядок их
% создания; если топологии нет — опишите ведение ваших основных сущностей.
\begin{hseExample}
Оператор может выполнять проверку и ведение топологии:

\begin{enumerate}[label=\arabic*), leftmargin=1.25cm]
    \item создать или выбрать \hseFill{родительская сущность топологии};
    \item создать \hseFill{дочерние сущности} и зарегистрировать \hseFill{конечные узлы или элементы};
    \item убедиться, что топология доступна для просмотра через API.
\end{enumerate}

Ожидаемый результат: сущности топологии связаны и возвращаются в корректной иерархии; ошибки валидации возвращаются при некорректных идентификаторах и атрибутах.
\end{hseExample}

\subsection{Управление сервисным каталогом}

\begin{hseExample}
При работе с сервисным каталогом оператор проверяет:

% Подсказка: скорректируйте перечень проверок под ваш каталог сервисов.
\begin{itemize}[leftmargin=1.25cm]
    \item создание \hseFill{типы сервисов или объектов каталога};
    \item фиксацию зависимостей между сервисами;
    \item просмотр графа зависимостей или списка связанных сервисов;
    \item корректное отклонение некорректных зависимостей.
\end{itemize}

Ожидаемый результат: сервисы и зависимости сохраняются в системе, а некорректные зависимости отклоняются с управляемой ошибкой.
\end{hseExample}

\subsection{Контроль пользовательских сценариев}

\begin{hseExample}
Оператор может выполнять контрольные пользовательские сценарии, не изменяя исходный код:

% Подсказка: замените перечень на контрольные сценарии вашей системы.
\begin{enumerate}[label=\arabic*), leftmargin=1.25cm]
    \item регистрация пользователя;
    \item подтверждение адреса электронной почты или номера телефона;
    \item вход с паролем и дополнительной проверкой;
    \item восстановление доступа;
    \item проверка активной сессии;
    \item выход из системы или отзыв токена.
\end{enumerate}

При каждом сценарии оператор проверяет HTTP-код ответа, тело ответа, отправку уведомления при сценариях с одноразовыми кодами, состояние сессии и отсутствие непредусмотренных ошибок в \hseFill{средство сбора ошибок}.
\end{hseExample}

\subsection{Просмотр журналов, метрик и трассировок}

\begin{hseExample}
Для первичной диагностики оператор использует:

\begin{itemize}[leftmargin=1.25cm]
    \item \hseFill{средство мониторинга} для просмотра метрик сервисов и инфраструктуры;
    \item \hseFill{средство сбора журналов} для поиска журналов по идентификатору запроса, имени сервиса или уровню ошибки;
    \item \hseFill{средство трассировки} для просмотра распределенной трассировки запроса;
    \item \hseFill{средство сбора ошибок} для анализа групп ошибок приложения.
\end{itemize}

При анализе инцидента оператор фиксирует время события, имя сервиса, идентификатор запроса, пользователя, код ошибки и ссылку на соответствующую панель или трассировку.
\end{hseExample}

\subsection{Завершение работы}

\begin{hseExample}
После выполнения операторских действий необходимо:

\begin{enumerate}[label=\arabic*), leftmargin=1.25cm]
    \item сохранить результаты проверки или ссылку на артефакты испытаний;
    \item удалить временные токены и файлы с чувствительными данными;
    \item остановить локальные контейнеры, если они больше не требуются;
    \item передать выявленные ошибки ответственному инженеру с приложением журналов, метрик или трассировок;
    \item зафиксировать состояние стенда, если проверка выполнялась в приемочном контуре.
\end{enumerate}
\end{hseExample}

\clearpage

% ============================================================
% 4. СООБЩЕНИЯ ОПЕРАТОРУ
% ============================================================
\section{СООБЩЕНИЯ ОПЕРАТОРУ}

\subsection{Общие правила обработки сообщений}

% Подсказка: скорректируйте перечень видов сообщений под вашу систему
% (HTTP-ответы, записи журналов, статусы задач, уведомления).
\begin{hseExample}
Сообщения оператору выдаются в виде HTTP-кодов и JSON-ответов API, записей журналов, уведомлений в средствах наблюдаемости и статусов задач CI/CD.

При получении ошибки оператор должен:

\begin{enumerate}[label=\arabic*), leftmargin=1.25cm]
    \item определить источник сообщения: API, сервис, инфраструктура, CI/CD или система наблюдаемости;
    \item зафиксировать код ответа, текст сообщения, время, пользователя и идентификатор запроса;
    \item повторить операцию только если сообщение допускает повтор;
    \item при системной ошибке проверить журналы, метрики и состояние зависимостей;
    \item передать информацию ответственному инженеру, если ошибка не устраняется операторскими действиями.
\end{enumerate}
\end{hseExample}

\subsection{Сообщения API и действия оператора}

% Подсказка: типовые HTTP-коды универсальны — скорректируйте формулировки
% под вашу систему и добавьте специфичные для проекта сообщения.
\begin{longtable}{|>{\raggedright\arraybackslash}p{0.22\textwidth}|>{\raggedright\arraybackslash}p{0.35\textwidth}|>{\raggedright\arraybackslash}p{0.33\textwidth}|}
\caption{Типовые сообщения программы и действия оператора}
\label{tab:ro_messages}\\
\hline
\textbf{Сообщение / код} & \textbf{Содержание сообщения} & \textbf{Действия оператора} \\
\hline
\endfirsthead
\caption*{Продолжение таблицы \thetable}\\
\hline
\textbf{Сообщение / код} & \textbf{Содержание сообщения} & \textbf{Действия оператора} \\
\hline
\endhead
\texttt{200 OK} & Операция выполнена успешно, данные возвращены в ответе & Проверить ожидаемое состояние объекта и, если сценарий подключен к аудиту, наличие события аудита \\
\hline
\texttt{201 Created} & Объект создан: \hseFill{перечислите типы объектов} & Сохранить идентификатор объекта и проверить его доступность через запрос чтения \\
\hline
\texttt{204 No Content} & Операция выполнена без тела ответа & Проверить результат повторным запросом или по событию аудита \\
\hline
\texttt{400 Bad Request} & Некорректный запрос: ошибка формата, неверный параметр, отсутствующее поле & Проверить тело запроса, обязательные поля, типы данных и формат идентификаторов \\
\hline
\texttt{401 Unauthorized} & Пользователь не аутентифицирован или токен недействителен & Повторить вход, проверить срок действия токена и корректность заголовка \texttt{Authorization} \\
\hline
\texttt{403 Forbidden} & Пользователь не имеет прав на выполнение операции & Проверить права доступа пользователя и корректность выбранного контура \\
\hline
\texttt{404 Not Found} & Запрошенная сущность не найдена & Проверить идентификатор, окружение и наличие объекта в списке \\
\hline
\texttt{409 Conflict} & Конфликт состояния: объект уже существует, операция нарушает доменное правило & Проверить уникальность имени или идентификатора; повторять запрос без изменения параметров не следует \\
\hline
\texttt{422 Validation Error} & Запрос не соответствует схеме API & Сверить запрос с OpenAPI-описанием и исправить структуру JSON \\
\hline
\texttt{429 Too Many Requests} & Превышено ограничение частоты запросов или отправки уведомлений & Дождаться истечения интервала ограничения, не запускать массовый повтор без согласования \\
\hline
\texttt{500 Internal Server Error} & Непредусмотренная ошибка сервиса & Зафиксировать идентификатор запроса, проверить журналы сервиса и \hseFill{средство сбора ошибок}, передать инцидент инженеру \\
\hline
\texttt{503 Service Unavailable} & Недоступна зависимость: \hseFill{СУБД, кеш, брокер сообщений} & Проверить состояние зависимостей, готовность сервисов и журналы инфраструктуры \\
\hline
\texttt{Health check not ready} & Сервис запущен, но еще не готов принимать трафик & Дождаться готовности или проверить миграции и подключения к зависимостям \\
\hline
\texttt{Pipeline failed} & Задание CI/CD завершилось ошибкой & Открыть журнал задания, определить этап сбоя и передать ссылку ответственному инженеру \\
\hline
\texttt{Alert firing} & Сработало правило мониторинга & Проверить панель мониторинга, журналы и трассировки; зафиксировать время начала и затронутые сервисы \\
\hline
\end{longtable}

\subsection{Сообщения уведомлений}

\begin{hseExample}
В сценариях регистрации, восстановления доступа и подтверждения операций программа может формировать уведомления с одноразовыми кодами через \hseFill{сервис или канал уведомлений}. Оператор должен учитывать следующие случаи:

\begin{itemize}[leftmargin=1.25cm]
    \item сообщение отправлено успешно: в ответе API указан идентификатор сообщения или подтверждение отправки;
    \item адрес назначения некорректен: запрос отклоняется ошибкой валидации;
    \item превышен лимит отправки: оператор ожидает истечения окна ограничения;
    \item провайдер доставки недоступен: оператор проверяет журналы сервиса уведомлений и состояние внешней интеграции;
    \item код подтверждения истек: пользователь должен запросить новый код.
\end{itemize}
\end{hseExample}

\subsection{Сообщения аудита}

\begin{hseExample}
События аудита используются оператором для подтверждения сценариев безопасности. Оператор проверяет журнал событий через \hseFill{конечная точка или интерфейс аудита} и фиксирует следующие поля события:

\begin{itemize}[leftmargin=1.25cm]
    \item субъект действия: пользователь или сервисная учетная запись;
    \item тип действия;
    \item ресурс и идентификатор ресурса;
    \item результат выполнения;
    \item время события;
    \item контекст запроса при наличии.
\end{itemize}

Если ожидаемое событие безопасности отсутствует, оператор должен проверить, был ли запрос фактически выполнен, корректно ли выбран контур, не возникла ли ошибка публикации события и доступен ли сервис аудита.
\end{hseExample}

\clearpage

% Страницы источников печатаются без нижнего штампа (как в реальных
% комплектах) — только номер листа и код документа сверху.
\pagestyle{appendix}

% ============================================================
% СПИСОК ИСПОЛЬЗОВАННЫХ ИСТОЧНИКОВ
% ============================================================
\section*{СПИСОК ИСПОЛЬЗОВАННЫХ ИСТОЧНИКОВ}
\addcontentsline{toc}{section}{СПИСОК ИСПОЛЬЗОВАННЫХ ИСТОЧНИКОВ}

% Подсказка: добавьте ПЕРЕД ГОСТами источники по технологиям вашего
% проекта в алфавитном порядке, по образцу:
% \bibitem{grafana_docs}
% Grafana Documentation [Электронный ресурс]. Режим доступа: \url{https://grafana.com/docs/grafana/latest/}, свободный. (дата обращения: 14.03.2026).
\renewcommand{\refname}{}
\begin{thebibliography}{99}
\bibitem{gost19101}
ГОСТ 19.101-77: Виды программ и~программных документов. // Единая система программной документации.~--- М.: ИПК Издательство стандартов, 2001.

\bibitem{gost19103}
ГОСТ 19.103-77: Обозначения программ и~программных документов. // Единая система программной документации.~--- М.: ИПК Издательство стандартов, 2001.

\bibitem{gost19104}
ГОСТ 19.104-78: Основные надписи. // Единая система программной документации.~--- М.: ИПК Издательство стандартов, 2001.

\bibitem{gost19105}
ГОСТ 19.105-78: Общие требования к~программным документам. // Единая система программной документации.~--- М.: ИПК Издательство стандартов, 2001.

\bibitem{gost19106}
ГОСТ 19.106-78: Требования к~программным документам, выполненным печатным способом. // Единая система программной документации.~--- М.: ИПК Издательство стандартов, 2001.

\bibitem{gost19505}
ГОСТ 19.505-79: Руководство оператора. Требования к~содержанию и~оформлению. // Единая система программной документации.~--- М.: ИПК Издательство стандартов, 2001.

\bibitem{gost19604}
ГОСТ 19.604-78: Правила внесения изменений в~программные документы, выполненные печатным способом. // Единая система программной документации.~--- М.: ИПК Издательство стандартов, 2001.
\end{thebibliography}

\clearpage

% ============================================================
% ЛИСТ РЕГИСТРАЦИИ ИЗМЕНЕНИЙ
% ============================================================
\input{../common/change_log}

\end{document}
